\chapter{はじめに}

TODO: 本書の目的・対象読者・読み方

\section*{本書の表記}

本書には、次の 2 種類の枠が出てきます。
どちらも、\textbf{実際に手を動かすところ}です。
読むだけで済ませずに、自分のコンピュータで入力し、結果を確かめながら読み進めてください。

\subsection*{ターミナルの枠}

ウィンドウのような見た目の枠は、ターミナルで入力するコマンドと、その実行結果を表しています。
本文中でターミナルの枠が登場した場合は、実際にターミナルを開いたあと、コマンドを入力して、同じ結果になるかを確かめましょう。

\begin{terminal}[ターミナルの枠の例]
$ echo "Hello, robot!"
Hello, robot!
\end{terminal}

\begin{itemize}
  \item \cmd{\$} から始まる行は、一般ユーザで実行するコマンドです。\cmd{\$} は入力せず、その後ろの部分だけを入力して \key{Enter} を押します。
  \item \cmd{\#} から始まる行は、root 権限で実行するコマンドです。
  \item \cmd{\$} や \cmd{\#} が付いていない行は、コマンドの実行結果です（入力しません）。長い実行結果は、途中を \cmd{...} で省略しています。
  \item \cmd{<ファイル名>} のように山括弧で囲まれた部分は、自分の環境に合わせて置き換えてください。
\end{itemize}

\subsection*{ソースコードの枠}

行番号の付いた枠は、プログラムや設定ファイルの中身（ソースコード）を表しています。
エディタ（\ref{chap:editor}章）でファイルを作って同じ内容を書き、保存してから実行してみましょう。
左端の行番号は説明のためのもので、入力しません。

\begin{lstlisting}[style=python, caption={ソースコードの枠の例（hello.py）}]
name = "robot"
print(f"Hello, {name}!")
\end{lstlisting}

枠の上に書かれた \cmd{hello.py} のような名前は、保存するときのファイル名の例です。

\section*{動作環境}
\begin{table}[h]
  \centering
  \begin{tabular}{ll}
    \toprule
    項目 & バージョン \\
    \midrule
    実行環境 & VirtualBox の仮想マシン \\
    OS & Ubuntu 24.04 LTS \\
    ROS & ROS 2 Jazzy Jalisco \\
    言語 & Python 3.12 \\
    シミュレータ & Gazebo Harmonic \\
    \bottomrule
  \end{tabular}
\end{table}
